\chapter{Research Methodology}
\label{chap:met}
\chaptermark{Research Methodology}

This chapter describes the research methodology employed to evaluate Java garbage collector performance in Spring Boot web application workloads. Section~\ref{sec:research-design} outlines the overall research design. Section~\ref{sec:test-app} describes the benchmark application architecture. Section~\ref{sec:workload-design} details the workload scenario design. Section~\ref{sec:measurement} covers measurement instrumentation and data collection. Section~\ref{sec:statistical} addresses statistical analysis and validity considerations.

\section{Research Design}
\label{sec:research-design}

This study adopts a \textbf{controlled experimental design} in which the independent variable is the JVM garbage collector configuration and the dependent variables are application-level latency percentiles, throughput, and GC pause characteristics. All other variables --- application code, database schema, data volume, JVM version, and hardware platform --- are held constant across experimental runs.

The experimental approach was selected over simulation or analytical modelling because GC behavior is highly sensitive to runtime interactions between the JVM, the operating system scheduler, and the application's actual allocation patterns. These interactions cannot be reliably captured by analytical models and are poorly approximated by synthetic microbenchmarks \cite{blackburn2006dacapo}. Empirical measurement on a realistic application under controlled load is therefore the most valid methodology for the research questions addressed.

The study tests five garbage collectors available in JDK 21 LTS: Serial GC, Parallel GC, G1 GC, ZGC (standard and Generational), and Shenandoah GC. Four workload scenarios and three heap configurations are crossed in a factorial design, yielding a comprehensive coverage of the parameter space relevant to production Spring Boot deployments.

\section{Benchmark Application Architecture}
\label{sec:test-app}

\subsection{Technology Stack}
\label{subsec:tech-stack}

The benchmark application is implemented as a Spring Boot 3.2 REST service \cite{springboot2024}, chosen because Spring Boot represents the dominant framework for enterprise Java web development \cite{stackoverflow2025survey}. The persistence layer uses Spring Data JPA with Hibernate ORM 6.4 \cite{hibernate2024}, backed by PostgreSQL 16 \cite{postgresql2024} as the relational database. Connection pooling is provided by HikariCP 5.0 \cite{hikaricp2024} with a pool size of 20 connections, matching the maximum concurrent thread count in baseline test scenarios.

Application-level caching is implemented via Caffeine 3.x \cite{caffeine2024}, configured with a maximum size of 1,000 entries and a 300-second expiry. This represents a realistic two-level caching architecture: Caffeine provides the application-level (L1) cache, while Hibernate's first-level session cache provides request-scoped object identity guarantees.

The complete source code, JMeter test plans, GC log files, and Python analysis scripts are publicly available in the accompanying repository.\footnote{\url{https://github.com/Dema-koder/gc-benchmark}}

\subsection{Domain Model}
\label{subsec:domain-model}

The application models a simplified e-commerce domain comprising three JPA entities:

\begin{itemize}
    \item \textbf{User} --- represents a registered customer with email, name, and registration timestamp. Mapped to the \texttt{users} table with a sequence-based primary key (\texttt{allocationSize=50}).
    \item \textbf{Product} --- represents a catalogue item with name, price, stock quantity, and description. Mapped to the \texttt{products} table. Product entities are cached at the application level via Caffeine, representing long-lived objects that stress old generation collection.
    \item \textbf{Order} --- represents a purchase order linking a user to one or more products via a \texttt{@ManyToMany} join table (\texttt{order\_products}). Order creation involves multi-table writes within a single transaction, generating a characteristic burst of short-lived ORM proxy objects.
\end{itemize}

This domain model was designed to reproduce the allocation patterns characteristic of real Spring Boot applications: request-scoped object lifecycles for order creation (aligning with the generational hypothesis \cite{jones2011garbage, ungar1984generation}), intermediate-lived objects for JPA lazy-loading proxies, and long-lived objects for cached product entities.

\subsection{REST API Endpoints}
\label{subsec:endpoints}

Three REST endpoints of varying computational complexity are exposed:

\begin{enumerate}
    \item \textbf{\texttt{GET /api/products/\{id\}}} --- lightweight endpoint that retrieves a single product by primary key. On cache hit (Caffeine), no database I/O occurs; the response is assembled entirely from heap-resident objects. On cache miss, a single \texttt{SELECT} query is executed and the result is cached. This endpoint generates minimal object allocation per request ($\sim$5--10 short-lived objects for DTO mapping and JSON serialization).

    \item \textbf{\texttt{GET /api/orders/user/\{userId\}}} --- medium-weight endpoint executing a \texttt{JOIN FETCH} JPQL query across the \texttt{orders}, \texttt{order\_products}, and \texttt{products} tables. The result set is mapped to a list of \texttt{OrderDto} records. This endpoint generates moderate allocation: $\sim$20--50 objects per request depending on order count, including Hibernate result transformers, DTO instances, and JSON serializer buffers.

    \item \textbf{\texttt{POST /api/orders}} --- heavyweight endpoint creating a new order within a database transaction. The operation involves: user entity lookup (cache hit), product entity batch fetch, \texttt{Order} entity instantiation, \texttt{order\_products} join table insertion, and \texttt{OrderDto} mapping. This endpoint generates the highest allocation density ($\sim$50--200 objects per request) including JPA entity instances, Hibernate proxy wrappers, transaction management objects, and JDBC statement wrappers.
\end{enumerate}

\subsection{Database Configuration}
\label{subsec:db-config}

PostgreSQL 16 was used as the database server, running on the same machine as the application to eliminate network latency variability. The database schema includes the following indices to ensure query performance is not a confounding variable: \texttt{idx\_orders\_user\_id} on \texttt{orders(user\_id)}, and \texttt{idx\_orders\_status} on \texttt{orders(status)}.

The test dataset comprises 500 users, 200 products, and 2,000 orders (generated with fixed random seed 42 for reproducibility). Sequence generators use \texttt{allocationSize=50} to batch primary key allocation and reduce round-trips to the PostgreSQL sequence, reflecting best-practice Hibernate configuration \cite{hibernate2024}.

\section{Workload Scenario Design}
\label{sec:workload-design}

\subsection{Load Generation Tool}
\label{subsec:jmeter}

Apache JMeter 5.6.3 \cite{jmeter2024} was selected as the load generation tool due to its support for HTTP keep-alive connection reuse, configurable thread groups, and per-request timing measurement. JMeter was executed in CLI (non-GUI) mode to eliminate GUI rendering overhead from measurements. All test plans were authored in JMX format and version-controlled alongside the application source.

A 100\,ms constant timer was inserted between request iterations within each thread to prevent TCP port exhaustion on the macOS loopback interface, a platform-specific constraint arising from the operating system's default ephemeral port range and TIME\_WAIT socket state duration. This think time also serves to simulate realistic inter-request intervals for a single user session.

\subsection{Workload Scenarios}
\label{subsec:scenarios}

Four scenarios were designed to capture the diversity of production web application traffic patterns:

\begin{enumerate}
    \item \textbf{Read-Heavy} (baseline): 20 concurrent threads executing a mixed request sequence --- \texttt{GET /api/products/\{id\}}, \texttt{GET /api/orders/user/\{userId\}}, \texttt{GET /api/products} --- with a 20\% probability of appending a \texttt{POST /api/orders} request. Duration: 300\,s. This scenario represents a typical content-serving or catalogue-browsing service with predominantly read traffic.

    \item \textbf{Write-Heavy}: 20 concurrent threads with inverted composition --- 80\% \texttt{POST /api/orders}, 20\% \texttt{GET /api/products/\{id\}}. Duration: 300\,s. This scenario represents transaction-intensive services (e.g., order processing, event ingestion) where write operations dominate and object allocation rates are substantially higher.

    \item \textbf{High Concurrency}: 50 concurrent threads executing the read-heavy request mix with 100\,ms think time. Duration: 300\,s. This scenario tests GC behavior under increased thread contention for database connections and application resources, representing services operating near their concurrency limit.

    \item \textbf{Burst}: A three-phase load profile implemented using three independent JMeter thread groups with scheduled delays:
    \begin{itemize}
        \item Phase 1 (Normal): 20 threads, 0--90\,s
        \item Phase 2 (Burst): 60 threads, 90--150\,s (3$\times$ normal concurrency)
        \item Phase 3 (Recovery): 20 threads, 150--240\,s
    \end{itemize}
    This scenario simulates realistic traffic spikes such as flash sales, batch job triggers, or upstream service recovery events.
\end{enumerate}

User and product IDs in all requests are selected uniformly at random from the test dataset range using JMeter's \texttt{RandomVariableConfig} element, ensuring no systematic cache warm-up bias toward specific entities.

\section{GC Configuration and Measurement}
\label{sec:measurement}

\subsection{JVM Configuration}
\label{subsec:jvm-config}

Each experimental run launches a fresh JVM instance with the following base configuration:

\begin{itemize}
    \item JDK: Eclipse Temurin 21.0.5 (LTS)
    \item Heap: \texttt{-Xms<50\% of max>} and \texttt{-Xmx<target>}, where target $\in \{512\,\text{MB}, 2\,\text{GB}, 4\,\text{GB}\}$
    \item GC logging: \texttt{-Xlog:gc*:file=<path>:time,uptime,level,tags}
    \item All other JVM flags: JVM ergonomic defaults
\end{itemize}

The collector-specific flag is the sole variable between runs:

\begin{center}
\begin{tabular}{ll}
\toprule
\textbf{Collector} & \textbf{Flag} \\
\midrule
Serial GC      & \texttt{-XX:+UseSerialGC} \\
Parallel GC    & \texttt{-XX:+UseParallelGC} \\
G1 GC          & \texttt{-XX:+UseG1GC} \\
ZGC            & \texttt{-XX:+UseZGC} \\
ZGC Gen.       & \texttt{-XX:+UseZGC -XX:+ZGenerational} \\
Shenandoah GC  & \texttt{-XX:+UseShenandoahGC} \\
\bottomrule
\end{tabular}
\end{center}

No GC-specific tuning parameters (pause time targets, region sizes, thread counts) were applied beyond the collector selection flag. This decision reflects the realistic scenario of a team adopting a new collector with default configuration, and avoids the confound of optimizing parameters independently for each collector.

\subsection{Application Warm-Up}
\label{subsec:warmup}

Each JVM instance is given a 20-second warm-up period after startup before JMeter begins sending requests. During this period, the Spring Boot application context initialises, Hibernate validates the database schema, HikariCP establishes the connection pool, and the JVM loads and JIT-compiles the most frequently executed code paths. The DataLoader component populates the test dataset only on the first run; subsequent runs detect existing data and skip population.

\subsection{GC Log Analysis}
\label{subsec:gc-analysis}

GC pause events are extracted from unified GC logs using collector-specific regular expression patterns targeting stop-the-world pause events:

\begin{itemize}
    \item \textbf{G1/Parallel/Serial}: \texttt{Pause (Young|Mixed|Full).*(\textbackslash d+[,\.]\textbackslash d+)ms}
    \item \textbf{ZGC}: \texttt{Pause (Mark Start|Mark End|Relocate Start).*(\textbackslash d+[,\.]\textbackslash d+)ms}
    \item \textbf{Shenandoah}: \texttt{Pause (Init Mark|Final Mark|Init Update Refs|Final Update Refs).*(\textbackslash d+[,\.]\textbackslash d+)ms}
\end{itemize}

Decimal separator normalisation (comma to period) was applied to handle the Russian locale \texttt{ru\_RU} active on the test machine, which causes the JVM to format floating-point numbers with comma separators in log output.

\subsection{Latency Metrics}
\label{subsec:metrics}

JMeter records a timestamp and elapsed time for each HTTP sample. Post-experiment analysis computes the following metrics from the raw CSV output using Python (pandas, numpy):

\begin{itemize}
    \item \textbf{Throughput}: successful requests per second over the measurement window
    \item \textbf{Average latency}: arithmetic mean of elapsed times
    \item \textbf{Percentile latencies}: p95, p99, p999 computed as \texttt{numpy.quantile} at the corresponding quantile
    \item \textbf{Maximum latency}: single largest elapsed time observed
    \item \textbf{GC pause statistics}: count, mean, p99, and maximum pause duration extracted from GC logs
\end{itemize}

For the burst scenario, per-30-second windows are computed to enable phase-level analysis of latency behaviour during normal, burst, and recovery phases.

\section{Statistical Validity Considerations}
\label{sec:statistical}

\subsection{Measurement Precision}
\label{subsec:precision}

JMeter elapsed time is measured in milliseconds using \texttt{System.currentTimeMillis()}, providing 1\,ms resolution. At the observed latency range (2--30\,ms), this represents 3--50\% resolution for individual samples but is negligible for percentile statistics computed over 50,000+ samples per run. GC pause durations from unified logging have sub-millisecond precision (three decimal places) and are not subject to measurement quantisation.

\subsection{Single-Run Design}
\label{subsec:single-run}

Each GC-scenario-heap combination was measured in a single 300-second run rather than multiple repeated runs. This decision was driven by experimental scope constraints: with 5 collectors $\times$ 4 scenarios $\times$ 3 heap sizes, a full factorial design with 3 repetitions would require 180 five-minute runs ($\sim$15 hours of continuous testing). The 300-second window provides approximately 50,000--100,000 individual samples per run, yielding stable percentile estimates (standard error of p99 estimate $<$0.1\,ms for $n > 10,000$).

The primary threat from the single-run design is run-to-run variance due to OS scheduler jitter and JIT compilation non-determinism. This threat is mitigated by the large sample count and the use of percentile statistics, which are more robust to individual outliers than mean-based statistics.

\subsection{Experimental Isolation}
\label{subsec:isolation}

All experimental components --- the JVM under test, JMeter, and PostgreSQL --- ran on the same physical machine (Apple M3 Pro Max, 11 cores, 18\,GB unified memory, macOS 15.6). This shared-machine design introduces inter-process interference as a potential confound, particularly at the 4\,GB heap configuration where memory pressure from all three processes is significant. The implications of this design choice for result validity are discussed in Section~\ref{sec:threats}.

\section{Chapter Summary}
\label{sec:met-summary}

This chapter described the controlled experimental methodology used to evaluate GC performance. The benchmark application reproduces key allocation patterns of production Spring Boot services: request-scoped object lifecycles, ORM-generated object graphs, and multi-level caching. Four workload scenarios capture the diversity of real web application traffic. GC unified logging provides sub-millisecond precision pause measurements, and JMeter provides high-sample-count latency percentile data. The methodology directly addresses the application domain gap, metrics gap, and modern collector gap identified in the literature review (Section~\ref{sec:research-gaps}).
